\scanpage{92}
Anneaux de classes de cycles\\
sur une variété projective non singulière\\
sur un corps $k$ alg. clos

\diagram{\begin{tikzcd}[ampersand replacement=\&,
 column sep=1.2cm,row sep=1.25cm,
 execute at end picture={
  \draw ([xshift=-0.5cm,yshift=0.45cm]anum.north west) rectangle ([xshift=0.45cm,yshift=-0.45cm]az.south east);
  \draw ([xshift=-0.45cm,yshift=0.4cm]apic.north west) rectangle ([xshift=1.8cm,yshift=-0.45cm]alin.south east);
  \draw[->] (discrete.west) -- ([xshift=0.45cm,yshift=0.25cm]aq.north east);
  \draw[->] (continuous.west) -- ([xshift=1.8cm]alin.east);
  \draw ([xshift=-0.35cm]anum.south west) .. controls ([xshift=-0.65cm]anum.south west) and ([xshift=-0.3cm,yshift=0.4cm]atau.west) .. ([xshift=-0.65cm,yshift=0.6cm]atau.west)
        .. controls ([xshift=-0.3cm,yshift=0.3cm]atau.west) and ([xshift=-0.65cm]atau.west) .. (atau.west);
  \node[anchor=east,align=left] at ([xshift=-0.7cm,yshift=0.6cm]atau.west) {groupe libre\\de type fini?};
  \draw (atau.south west) .. controls ([xshift=-0.55cm]atau.south west) and ([xshift=-0.3cm,yshift=0.5cm]aalg.west) .. ([xshift=-0.65cm,yshift=0.7cm]aalg.west)
        .. controls ([xshift=-0.3cm,yshift=0.3cm]aalg.west) and ([xshift=-0.55cm]aalg.west) .. (aalg.west);
  \node[anchor=east] at ([xshift=-0.7cm,yshift=0.7cm]aalg.west) {Groupe fini?};
  \draw (aalg.south west) .. controls ([xshift=-0.4cm]aalg.south west) and ([xshift=-0.2cm,yshift=0.5cm]apic.west) .. ([xshift=-0.55cm,yshift=0.7cm]apic.west)
        .. controls ([xshift=-0.2cm,yshift=0.3cm]apic.west) and ([xshift=-0.4cm]apic.west) .. (apic.west);
  \node[anchor=east] at ([xshift=-0.6cm,yshift=0.7cm]apic.west) {V.A.?};
  \draw ([xshift=-0.3cm,yshift=0.3cm]oldq.north west)--([xshift=0.35cm,yshift=-0.4cm]oldz.south east);
  \draw ([xshift=0.2cm,yshift=0.35cm]oldq.north east)--([xshift=-0.3cm,yshift=-0.3cm]oldz.south west);
  \draw (car.west)--(az.north west);
  \draw ([yshift=-0.3cm]ck.south west) .. controls ([yshift=-0.5cm]ck.south west) and ([xshift=-0.15cm,yshift=-0.35cm]ck.south) .. ([yshift=-0.55cm]ck.south)
        .. controls ([xshift=0.15cm,yshift=-0.35cm]ck.south) and ([yshift=-0.5cm]ck.south east) .. ([yshift=-0.3cm]ck.south east);
  \node[anchor=north,align=left] at ([yshift=-0.65cm]ck.south) {flèches\\qui sont des\\isom mod torsion};
  \draw ([xshift=-0.25cm,yshift=-0.45cm]alin.south west) .. controls ([yshift=-0.65cm]alin.south west) and ([xshift=0.25cm,yshift=-0.45cm]alin.south) .. ([xshift=0.45cm,yshift=-0.7cm]alin.south)
        .. controls ([xshift=0.6cm,yshift=-0.45cm]alin.south) and ([xshift=1.65cm,yshift=-0.65cm]alin.south) .. ([xshift=1.65cm,yshift=-0.45cm]alin.south);
  \node[anchor=north,align=left] at ([xshift=0.7cm,yshift=-0.8cm]alin.south) {flèches toutes\\surjectives};
 }]
 \& |[alias=anum]| A_{\mathrm{num}} \&\&\& |[alias=discrete]|\text{théorie discrète.}\\
 \& |[alias=atau]| A_\tau \arrow[u,"{?\ \unsure{\wr}}"']\arrow[r,"{\overset{!}{\sim}}"]
 \& |[alias=aq]|A_{Q_\ell\text{-hom}} \arrow[r]
 \& |[alias=oldq]|\deleted{\unsure{A_{Q_\ell\text{-hom}}}}
 \& |[alias=car,align=left]|\ensuremath{\substack{\text{\textcircled{car. 0}}\\\text{\unsure{en remplaçant}}\\\text{\unsure{théorie sing. entière.}}}}\\
 \& |[alias=aalg]| A_{\mathrm{alg}}\arrow[u]\arrow[r,"{\overset{?}{\sim}\ \textcircled{\unsure{t0|th0}}}"]
 \& |[alias=az]| A_{\unsure{\wedge Z_\ell|\widehat Z_\ell}\text{-hom}}\arrow[u]\arrow[r]
 \& |[alias=oldz]|\deleted{\unsure{A_{Z_\ell\text{-hom}}}}\arrow[u]\\
 |[alias=gtop]| G_{\mathrm{top}}K\arrow[uur,bend left=55]\arrow[urr,dashed,"?",bend right=10]
 \& |[alias=apic]| A_{\mathrm{pic}}\arrow[u]\\
 \& A_{\mathrm{alb}}\arrow[u,"{\unsure{\wr}\ ?}"]\\
 G_{\mathrm{alg}}K\arrow[uu]
 \& |[alias=alin]| A_{\mathrm{lin}}\arrow[u,"{\unsure{\wr}\ ?}"]\arrow[uul,"{\text{surj}}"']
 \&\& |[alias=continuous]|\text{théorie continue}\\
 |[alias=ck]| CK\arrow[u,"{\text{surj}}"']\arrow[ur]
\end{tikzcd}}

\margin{Autres équivalences : l'équivalence \unsure{régulière|rigide} $r$ et l'équivalence de Weil $\unsure{wh|w_h}$ et $W$ définis par Liebermann \ldots\ldots}

\margin{91}
